\subsection{Strict cocycles as integral tangent functionals}
\label{ct:section}

Let \(O\) be the integers in a finite extension of \(\mathbb Q_\ell\),
\(0\ne a\in O\), \(M=O/aO\), and \(\Gamma=G_{\mathbb Q}\).
Write \(O\ltimes M=O\oplus\epsilon M\), with \(\epsilon^2=0\).
This is a complete local \(O\)-algebra. Its quotients
\((O/\pi^h)\ltimes M\), \(h\ge v_\pi(a)\), are Artin rings and their
inverse limit is \(O\ltimes M\). For a continuous augmented
\(O\)-algebra \(B\xrightarrow{x}O\), a continuous map to this
extension over \(x\) has the unique form \(x+\epsilon d\), with
\(d(fg)=x(f)d(g)+x(g)d(f)\).

Use one-based indices for matrix units. The two diagonal points
\(r_0\), their infinitesimal directions \(X\), and the weights \(j\)
of their nonzero entries are as follows:
\[
\begin{array}{c|c|c|c}
 n&\text{matrix of }r_0& X& j\\ \hline
 3&\operatorname{diag}(\Psi,1,\Psi^{-1})&E_{21}-E_{32}&-1\\
 4&\operatorname{diag}(\alpha\Psi,\Psi,\alpha,1)&E_{31}+E_{42}&(0,-1).
\end{array}
\]
In rank three one appends the sign character \(\eta\) to the matrix,
in the semidirect product in which the negative sign acts by
\(g\mapsto Jg^{-t}J\), \(J\) the antidiagonal matrix of ones. Its
differential \(Y\mapsto-JY^tJ\) fixes \(X\).
In rank four the symplectic adjoint sends \(X\) to \(-X\);
the diagonal point has similitude \(\alpha\Psi\).
In both ranks \(\operatorname{Ad}(r_0(g))X=\theta(g)X\) with
\(\theta=\Psi^{-1}\). The rank-four bidegrees of the four basis
vectors are \((1,1),(0,1),(1,0),(0,0)\).

\begin{proposition}[The square-zero representation]\label{ct:point}
Let \(z:\Gamma\to M\) be a continuous cocycle for \(\theta\).
Then \(r_z(g)=(1+\epsilon z(g)X)r_0(g)\) is a continuous
representation over \(O\ltimes M\) with the prescribed sign or
similitude. If on a subgroup \(D\) one has
\(z(g)=(\theta(g)-1)m\), the restriction is the reduction modulo \(a\)
in the second coordinate of an integral dual-number conjugate of
\(r_0|_D\).
\end{proposition}
\begin{proof}
Multiplication adds the two square-zero coefficients as
\(z(g)+\theta(g)z(h)=z(gh)\). In rank four the linear terms in the
similitude cancel because the symplectic adjoint of \(X\) is \(-X\).
Choose a lift \(\widetilde m\in O\) of \(m\). Conjugation by
\(1-\epsilon\widetilde mX\) changes the left infinitesimal coefficient
of \(r_0(g)\) to \(\widetilde m(\theta(g)-1)X\).
The matrix entries are continuous in the adic topology since \(r_0\)
is continuous and \(z\) is locally constant.
\end{proof}

For rank three take the normalized representative of each class in
\(H[b]\), where \(H\) is the strict group and \(0\ne b\in O\). At the
marked element \(g_1\) one has \(\Psi(g_1)^2-1\in O^\times\);
subtracting a unique coboundary makes \(z(g_1)=0\). The normalized
representative has values in \((\operatorname{Frac}(O)/O)[b]\simeq O/bO\).
Here the local group at \(\ell\) contains \(s\) with
\(\Psi(s)^2-1\in O^\times\), so its coboundary witness is also killed
by \(b\). Take \(a=b\). At complex conjugation \(c\), the identities
\(\theta(c)=1,\ c^2=1\) give \(2z(c)=0\), hence \(z(c)=0\) since
\(\ell\ne2\). The cocycle identity then gives \(z(cg_1^{-1}c)=0\).
Consequently the matrices at the three elements entering the two
normalization equations agree with those of the base point.
On the prescribed inertia elements away from \(\ell\) where the
residual representation is trivial, the datum requires \(r_0(g)=1\).
There \(\theta(g)=1\), so local strictness gives \(z(g)=0\).
\src{src/Development/AutomorphicArithmetic/AutomorphicVarieties/UnitaryThreeGroups/ShaInjGen.lean}{145}{normalized rank-three cocycles and their root entries}.

In rank three the local condition at \(\ell\) specifies particular
components of a closed locus. Its first-order verification can be made
directly. Let \(U\) represent unrestricted framed local lifts with the
prescribed sign. Put \(H_L=\ker\eta\cap\ker\tau\), where
\(\Psi|_{I_\ell}=\varepsilon_\ell^e\tau\), \(e\ge1\), with \(\tau\)
of finite order. Let \(I_{\rm ord}\subset U\) be the ideal of the
inertia characteristic-polynomial identities and the ordered products
of three shifted matrices on \(H_L\cap I_\ell\), with successive
characters \(\Psi,1,\Psi^{-1}\). Write \(\Delta=\ell d_1d_2\) for the
product of the two consecutive Frobenius-gap equations used to saturate
this ideal. Let \(I_t\) be the intersection of the minimal primes of
\(\sqrt{I_{\rm ord}:\Delta^\infty}\) contained in the kernel of the
diagonal point. The framed universal property makes this point unique.
\src{src/Development/AutomorphicArithmetic/AutomorphicVarieties/PadicHodgeTheory/IsInertiaInfiniteCyclotomic.lean}{133}{the selected ordinary components}.

\begin{lemma}\label{ct:local}
For \(Y\in M_3(O)\), the local framed representation
\((1+\epsilon Y)r_0(1+\epsilon Y)^{-1}\) over \(O[\epsilon]\),
\(\epsilon^2=0\), kills \(I_t\).
\end{lemma}
\begin{proof}
Let \(d:G_{\mathbb Q_\ell}\to\widehat{\mathbb Z}\to\mathbb Z_\ell\)
be unramified degree, normalized on arithmetic Frobenius. The character
\(u_T(g)=(1+T)^{d(g)}\) is continuous for the \((\pi,T)\)-adic
topology by the binomial expansion. Over the domain \(O[[T,V]]\)
use the diagonal point for \(\Psi u_T\), with the same sign as before,
conjugated by \(1+VY\). The antidiagonal inverse-transpose involution
fixes every matrix \(\operatorname{diag}(b,1,b^{-1})\), so this
does define a representation of the full local group. On
\(H_L\cap I_\ell\) the diagonal characters are unchanged. The ordered
product of the three shifted diagonal matrices is zero since each
diagonal entry has its own vanishing factor. Conjugation preserves this
identity and the three characteristic coefficients.

At the distinguished inertia element \(s_0\in H_L\) the three diagonal
values \(\lambda_i\) are distinct: \(\varepsilon_\ell(s_0)\) has
infinite order. Set \(c_i=\prod_{j\ne i}(\lambda_i-\lambda_j)\).
The product of the other two shifted matrices is \(c_iE_{ii}\) in
the diagonal basis. At the specified \(\phi\in H_L\) of nonzero
unramified degree, put \(b=\Psi(\phi)(1+T)^{d(\phi)}\). Its three
diagonal values \(b_i\) are \(b,1,b^{-1}\). The trace defining the
\(i\)-th projected Frobenius value is \(c_i b_i\), and hence
\[
 d_i=c_i c_{i+1}\bigl(b_i-\varepsilon_\ell(\phi)b_{i+1}\bigr)
 \qquad(i=1,2).
\]
These are invariant under the displayed conjugation since all the
elements in the traces have positive sign. Both are nonzero: \(b\)
has nonzero linear coefficient. Thus the map \(U\to O[[T,V]]\)
kills \(I_{\rm ord}\) and its saturation and radical. Its prime kernel
contains a minimal prime \(P\) of that radical. Evaluation at
\(T=V=0\) is the diagonal point, so \(P\) occurs in the intersection
defining \(I_t\). In particular the map kills \(I_t\). Specializing
\(T=0,\ V=\epsilon\) proves the assertion, including when a gap
vanishes at the original diagonal point.
\end{proof}

In rank four use the common denominator for the local witnesses as
in Lemma~\ref{newcoh:cotangent} and the paragraph following it.
At \(\ell\) the indicated ordinary equations prescribe the
characteristic polynomial at \(s_0\), the vanishing of the forbidden
entries between its polynomial projectors, and the prescribed diagonal
inertia characters between equal projectors. These hold on the
diagonal representation and survive symplectic conjugation, since a
polynomial in a matrix transforms by the same conjugation. At
\(p\in S_b=\{p\in S:p\ne\ell,\ \ell\mid\#\psi(I_p)\}\) the condition
is fixed similitude and triviality on \(I_p\cap\ker\psi\); the diagonal
datum has this property, and conjugation preserves it. At all of these
places Proposition~\ref{ct:point} supplies an integral conjugate using
the local witness. Away from \(T=\{\ell\}\cup S_b\), the datum has
\(r_0(g)=1\) on inertia elements with trivial residual action, and the
strict cocycle vanishes there. Thus the local lifts and global \(r_z\)
agree after reduction, with every comparison frame equal to the identity.
The representing framed global ring consequently gives a map into
\(O\ltimes M\); its restriction to the local coefficient ring lifts
to \(O[\epsilon]/(\epsilon^2)\). The same assertions hold for the
normalized rank-three ring by Lemma~\ref{ct:local} and the checks above.
\src{src/Development/AutomorphicArithmetic/AutomorphicVarieties/SymplecticGSp4/ExistsLocalFirstOrderLifts.lean}{188}{integral local conjugates}
and \src{src/Development/AutomorphicArithmetic/AutomorphicVarieties/SymplecticGSp4/ExistsQuotientDerivationOfShaCocycle.lean}{252}{the global framed point}.

\begin{lemma}[Integral lift on a relative presentation]\label{ct:lift}
Let \(L\) be an augmented complete local \(O\)-algebra with weighted
coordinates at the augmentation, and \(P=L[[t_1,\ldots,t_s]]\) with
homogeneous variables of augmentation zero. Suppose a map
\(P\to O\ltimes M\) over the augmentation restricts on \(L\) to the
reduction of an integral derivation of weight \(j\) in its second
coordinate. Suppose its second coordinate on \(t_i\) is zero when
\(\mathrm{wt}(t_i)\ne j\). Then its second coordinate lifts to an
integral derivation of weight \(j\) on \(P\).
\end{lemma}
\begin{proof}
Lift its value on each \(t_i\) to \(d_i\in O\), taking \(d_i=0\) in
all other weights. If \(D_L\) is the given local lift and \(x_L\) is
the augmentation of \(L\), the required map is
\[
 D\Bigl(\sum_\nu b_\nu t^\nu\Bigr)
     =D_L(b_0)+\sum_i x_L(b_{e_i})d_i .
\]
Comparing constant and linear coefficients proves Leibniz and
reduction. The augmentation is fixed by the grading, so this map
vanishes on the other homogeneous pieces.
\end{proof}

Apply the lemma on the weighted presentation before the global
relations. The presentation variables are shifted matrix entries or
comparison-frame entries. Derivatives of the frames vanish, and the
matrix derivatives are precisely the entries of \(z(g)Xr_0(g)\).
The rank-three presentation also uses the inverse-transpose block:
the derivative of inversion at a diagonal \(D\) is
\(Y\mapsto-D^{-1}YD^{-1}\); transposition and reversal of the indices
preserve the assigned block weight. Thus all these variables satisfy
the hypothesis on the weights. The local matrix-entry derivatives of
the integral conjugates have the same weight. More generally, suppose
maps from the local representing ring into integral dual numbers are
determined by the local matrix entries. For an integral torus unit
\(t\), precompose the local map by the weight action of \(t\), or
multiply its second coordinate by the character \(t^j\). The two maps
agree on all matrix entries and hence everywhere. On a homogeneous
element of weight \(i\) this gives \((t^i-t^j)d(f)=0\). For \(i\ne j\)
one can separate the two weights using \(1+\ell\) in one torus
coordinate. Since \(O\) is a domain this gives \(d(f)=0\). This
argument verifies the weight condition on the whole local derivation.
\src{src/Development/AutomorphicArithmetic/AutomorphicVarieties/UnitaryThreeGroups/ModelOfPres.lean}{302}{the relative rank-three presentation}.

In rank four the presentation is further cut by six homogeneous frame
or matrix entries per distinguished place, each with nonzero first
weight. The integral derivation of weight \((0,-1)\) kills them.
In the resulting quotient \(P'\), the model next divides by the kernel
of its characteristic-zero formal chart. The precise property needed
here is that for \(f\) in this kernel some \(0\ne c\in O\) satisfies
\(cf\in(\ker x_{P'})^2\). Then \(cD(f)=0\), so the integral derivation
descends through the chart kernel. This is the cotangent injectivity
condition of the chart after extension to \(\operatorname{Frac}(O)\).
The other normalization equations involve the frames at \(S_b\) and
off-diagonal entries of \(r(g_0)+\alpha(g_0)\Psi(g_0)r(g_0^{-1})\).
Their derivatives vanish modulo \(a\) because the frames are constant
and \(z(g_0)=z(g_0^{-1})=0\). All original global relations also vanish
modulo \(a\), since \(r_z\) defines a point of the framed ring.
\src{src/Development/AutomorphicArithmetic/AutomorphicVarieties/SymplecticGSp4/ExModel.lean}{496}{passage to the normalized weighted quotient}
and \src{src/Development/AutomorphicArithmetic/AutomorphicVarieties/SymplecticGSp4/ExistsSmoothQuotientByWeightedEquations.lean}{104}{the chart-kernel condition}.

On the sum of homogeneous pieces used in Lemma~\ref{newcoh:cotangent},
restrict the resulting derivation to weight \(-1\) (weight \((0,-1)\)
before taking the first-weight-zero algebra in rank four). It kills
the saturation of the square of the augmentation ideal because it
is integral, and its reduction kills the indicated equations.
The functional on \(Q_{-1}\) recovers \(z(g)\) from a homogeneous
lift of matrix entry \((2,1)\) or \((3,1)\), respectively, by dividing
by \(r_0(g)_{11}\in O^\times\). Distinct normalized cocycles therefore
give distinct functionals. In rank three take \(a=b\). In rank four
take \(a=b\prod_v(\theta(s_v)-1)\), as above, so that each local
witness is represented at this denominator. In either rank the same
\(Q_{-1}\) bounds the length of \(H[b]\) for every nonzero \(b\).
Every continuous cocycle into the discrete torsion module has finite
image by compactness, so \(H=\bigcup_hH[\pi^h]\). If \(Q_{-1}\)
has finite length, this increasing sequence stabilizes since a strict
inclusion would increase length. This proves finiteness and the
length bound for \(H\) in both ranks. The following sections construct the weighted presentations, count their
equations and prove cotangent injectivity at the rank-four cut.
Finiteness of the invariant quotient and rigidity at the remaining
global points are further requirements of Theorem~\ref{newcoh:orbitdegree}.
